% STATUS: verified
\chapter{The API}\label{ch:api}

% One route of a table: the methods and the path on a line of their own, across both columns;
% the row under it says who may call and what the route does.
\providecommand{\apiroute}[2]{\multicolumn{2}{@{}l@{}}{\route{#1}{#2}}\\*}

This chapter is the reference of the system's HTTP interface: every route, who may call it, what
it takes and what it refuses. It has five parts. \Cref{sec:api:conventions} gives the
conventions all routes share --- where a request goes, the token, the error envelope, paging,
limits. \Cref{sec:api:routes} lists the routes of the API (Django) area by area, one table for
each. \Cref{sec:api:assistant} gives the three public routes of the assistant service, which
carry the chat. \Cref{sec:api:internal} describes the internal doors the two services open to
each other and to nobody else. \Cref{sec:api:client} gathers what a program calling the API has
to know about work that outlasts a request, and about files. A \term{client} in this chapter is
any program that calls the API: the console, the assistant service acting for a user, or a
script.

The chapter says what each route does, not why: the reasons are in the chapter on the route's
subject, named beside each table. The file \file{urls.py} of the API is the table of contents
of the interface --- every route of the API is one line there --- and each views module opens
with a list of its routes and their bodies.


% ============================================================
\section{Conventions}\label{sec:api:conventions}
% ============================================================

\subsection{Three families of paths}\label{sec:api:families}

The system has one public address. Behind it stands the \term{ingress}, the Caddy server
\code{legalops-endpoint}, configured in \file{endpoint/Caddyfile}. It listens on port~80 ---
the only port the installation publishes --- and sorts every request by its path, the first
match winning (\cref{fig:api:routing}):

\begin{enumerate}
  \item \code{/api/assistant/*} goes to the assistant service (\code{legalops-assistant},
    port 3000). These are the chat's three routes (\cref{sec:api:assistant}).
  \item \code{/api/*} goes to the API (\code{legalops-django}, port 8000): every other route
    of this chapter that a client can reach.
  \item Everything else goes to the web server of the console and the public site
    (\code{legalops-vite}), which answers with files, never with data (\cref{ch:console}).
\end{enumerate}

One more path stands in front of these three: \code{/dbgate/*} leads to DBGate, the database
administration tool. Before the ingress passes such a request on, it asks the API whether the
caller is an administrator (\route{GET}{/api/auth/checkauth/admin}) and hands a refusal back as
it is (\cref{ch:ops}).

\begin{figure}[H]
\centering
\begin{tikzpicture}[node distance=4mm and 12mm]
  \node[svc, minimum width=44mm] (assistant) {assistant service\\{\scriptsize\code{legalops-assistant:3000}}};
  \node[svc, minimum width=44mm, below=of assistant] (django) {the API (Django)\\{\scriptsize\code{legalops-django:8000}}};
  \node[svc, minimum width=44mm, below=of django] (vite) {console and public site\\{\scriptsize\code{legalops-vite:80}}};
  \node[svc, minimum width=44mm, below=of vite] (dbgate) {DBGate\\{\scriptsize\code{legalops-dbgate:3000}}};
  \node[svc, minimum height=16mm, anchor=east] (in) at ([xshift=-36mm]$(django.west)!.5!(vite.west)$) {ingress\\{\scriptsize\code{legalops-endpoint}}\\{\scriptsize port 80}};
  \node[ext, left=8mm of in] (client) {browser\\or script};
  \node[note, below=3mm of in, text width=34mm] {before a \code{/dbgate/*} request it asks the API whether the caller is an administrator};
  \draw[flow] (client) -- (in);
  \draw[flow] (in.east) -- ++(4mm,0) |- (assistant.west) node[lbl, pos=.8] {\code{/api/assistant/*}};
  \draw[flow] (in.east) -- ++(4mm,0) |- (django.west) node[lbl, pos=.8] {\code{/api/*}};
  \draw[flow] (in.east) -- ++(4mm,0) |- (vite.west) node[lbl, pos=.8] {everything else};
  \draw[flow] (in.east) -- ++(4mm,0) |- (dbgate.west) node[lbl, pos=.8] {\code{/dbgate/*}};
  \begin{scope}[on background layer]
    \node[lane, fit=(assistant)(django)(vite)(dbgate)] (lane) {};
  \end{scope}
  \node[lanelabel, anchor=south west] at (lane.north west) {isolated network: no port of its own};
\end{tikzpicture}
\caption{Where a request goes. The ingress is the only service with a published port; it
chooses the service by the path alone. Nothing under \code{/internal/} or \code{/health} is
forwarded to the API or to the assistant service.}\label{fig:api:routing}
\end{figure}

Three things the ingress does concern every client. A request whose \code{Origin} header names
another site is answered 403 with an empty body before it reaches any service. The ingress
passes on the \code{X-Forwarded-For} header it received and copies it into \code{X-Real-IP};
that header is the address by which the API counts login attempts and writes audit rows
(\code{client_ip} in \file{common/http.py}). The value is whatever the incoming request
carried; without one, the API falls back to the address of the connection, which is the
ingress's own. And when the service behind a path is down or
restarting, the ingress itself answers 502 with an empty body. The security headers it adds to
every answer but its own 403 are described in \cref{ch:security}.

\subsection{Requests and answers}\label{sec:api:json}

The API speaks JSON only: a request body is a JSON object, and every answer --- an error
included --- is JSON written by \code{json_response} in \file{common/http.py}, with UTF-8 kept
literal. The exceptions are few and named in the tables: two uploads take
\code{multipart/form-data}, three exports answer a file, one route answers an SVG drawing, and
the chat answers a stream.

\begin{description}[style=unboxed, leftmargin=0pt]
  \item[Paths.] No path ends in a slash (\code{APPEND_SLASH} is off); the same path with one is
    a 404. An unknown path answers 404 \code{NOT_FOUND} and a method the route does not serve
    answers 405 \code{METHOD_NOT_ALLOWED} with an \code{Allow} header, both as JSON, never as an
    HTML page.
  \item[Field names.] camelCase (\code{displayName}, \code{companyId}). Two places differ: the
    chat's history rows keep the names of the chat runtime that reads them (\code{parent_id}),
    and the internal doors use snake\_case (\code{user_id}, \code{input_tokens}).
  \item[Values.] A datetime is ISO~8601 with its offset, as \code{2026-09-27T18:00:00+00:00};
    a date is \code{YYYY-MM-DD}; a cost is a number of US dollars.
  \item[Ids.] Every row a route creates gets a uuid4 string as its id. Prompts are addressed
    by their slug and a prompt version by its number. An id that arrives in a body or a query
    is accepted only as a non-empty string of at most 64 characters (\code{str_id}); anything
    else is treated as no id at all, so an object or a list never reaches a database lookup.
  \item[Transactions.] Every request runs in one database transaction (\code{ATOMIC_REQUESTS}
    in \file{settings.py}): a request that fails leaves nothing half-written. Four routes that
    wait on something slow stand outside it --- the health probe, the document upload, the
    prompt test run and the MCP connection test.
\end{description}

\subsection{The token and who may call}\label{sec:api:auth}

A client signs in with \route{POST}{/api/auth/login} and sends the token it receives in every
later request, as \code{Authorization: Bearer <token>}. The token is an opaque random string
(256 random bits; the database keeps only its SHA-256), valid for \code{SESSION_DAYS}~=~30 days;
an account keeps its \code{SESSIONS_PER_USER}~=~10 newest sessions, so an eleventh login ends
the oldest. The account is loaded again on every request, so a deactivated account is locked
out at once. There is no self-registration and no cookie authentication: the one cookie the
API sets, \code{compliance_gate}, is read only by the gate in front of DBGate. A request
without a valid token answers 401 \code{AUTH_REQUIRED}. The mechanism is in
\file{users/auth.py}; why it is built this way is argued in \cref{ch:security}.

There are two roles, \code{user} and \code{admin}. The tables of \cref{sec:api:routes} say who
may call a route with one of four words:

\begin{table}[H]
\caption{Who may call a route --- the words used in the route tables}\label{tab:api:who}
\begin{tabularx}{\textwidth}{@{}L{22mm}Y@{}}
\toprule
\hd{Word} & \hd{Means} \\
\midrule
public & No token is needed. Only the health probe and the login are public. \\
signed in & Any active account with a valid token (\code{@require_auth}). \\
owner & Any signed-in user, on the rows that are theirs; an administrator, on every row. Which
  rows are a user's own is the \term{tenancy rule}, below. \\
admin & Administrators only (\code{@require_role("admin")}, or a role check inside the view).
  A user gets 403 \code{FORBIDDEN}. \\
\bottomrule
\end{tabularx}
\end{table}

The tenancy rule is written once per kind of row, as a function that returns the rows a caller
may touch, and every route of that kind starts from it:

\begin{itemize}
  \item A \emph{company} belongs to the user who created it (\code{visible_companies} in
    \file{companies/api/company_views.py}). Its credentials, adopted standards, roles, reports,
    findings, actions and work products follow the company.
  \item A \emph{document} belongs to the user who uploaded it (\code{visible_documents}),
    whether or not it is attached to a company.
  \item A \emph{job} belongs to the user who caused it (\code{_visible} in
    \file{jobs/api/views.py}), and a \emph{model call} in the usage log to the user it is
    charged to.
  \item A \emph{chat thread} belongs to its user, and here an administrator is no exception
    (\code{visible_threads} in \file{chat/api/views.py}): nobody reads another person's
    conversations through the thread routes.
  \item The \emph{legal corpus} and the register of obligations are readable by everyone
    signed in. A user sees only the acts that are enabled; an administrator sees all. A
    company's own requirements are never listed by the obligation routes.
\end{itemize}

\paragraph{What a 404 means.} A row the caller may not see does not exist for that caller: the
route answers 404 \code{NOT_FOUND}, the same answer as for an id that was never issued, so the
API never confirms that someone else's company, document, report or thread exists. 403 is kept
for two cases where nothing is revealed by it: a user calling a route that only administrators
may call, and a user asking the usage routes for other accounts' figures.

\subsection{Errors}\label{sec:api:errors}

Every error of the API has one shape, the \term{error envelope} (\code{json_error} in
\file{common/http.py}):

\begin{codeblock}
{"error": "An analysis of this company is already running", "code": "ALREADY_RUNNING"}
\end{codeblock}

\code{error} is a sentence in English for a person; \code{code} is the stable name a client
branches on. The HTTP status carries the class of the failure. \Cref{tab:api:codes} lists every
code the API answers; the route tables name the ones particular to a route and leave out those
that any route can answer (\code{INVALID_BODY}, \code{AUTH_REQUIRED}, \code{FORBIDDEN},
\code{NOT_FOUND}, \code{METHOD_NOT_ALLOWED}).

\begin{table}[H]
\caption{The error codes of the API (Django), by HTTP status}\label{tab:api:codes}
\begin{tabularx}{\textwidth}{@{}L{12mm}Y@{}}
\toprule
\hd{Status} & \hd{Codes} \\
\midrule
400 & \code{INVALID_BODY} (no JSON object in the body), \code{MISSING_FIELDS},
  \code{INVALID_FIELD}, \code{INVALID_QUERY} (a query parameter), \code{EMPTY_UPDATE} (nothing
  to change), \code{MISSING_FILE}, \code{BAD_REQUEST} (a NUL character in the path or query),
  \code{INVALID_VALUE} (a value the database cannot take: out of range or malformed), \code{INVALID_EMAIL},
  \code{INVALID_ROLE}, \code{WEAK_PASSWORD}, \code{WRONG_PASSWORD}, \code{LAST_ADMIN},
  \code{SELF_LOCKOUT}, \code{UNKNOWN_MODEL}, \code{NOTHING_TO_ASSESS}, \code{NOTE_REQUIRED} \\
401 & \code{AUTH_REQUIRED} (no session, or a dead one), \code{INVALID_CREDENTIALS} (login) \\
403 & \code{FORBIDDEN} (the role), \code{ROLE_NOT_ALLOWED} (a chat message that is not the
  user's own), \code{INTERNAL_ONLY} (the internal door without its secret) \\
404 & \code{NOT_FOUND}; on the internal door also \code{PROMPT_NOT_CONFIGURED} \\
405 & \code{METHOD_NOT_ALLOWED}, with an \code{Allow} header \\
409 & The request is well formed but the row is not in a state for it:
  \code{ALREADY_RUNNING}, \code{STILL_RUNNING}, \code{NOT_CANCELLABLE}, \code{NOT_FINISHED},
  \code{NOT_WAITING}, \code{NOT_DRAFT}, \code{NOT_APPROVED}, \code{NOT_READY},
  \code{NEWER_VERSION}, \code{RULES_BROKEN}, \code{CLOSED_BY_REPORT}, \code{CORPUS_EMPTY},
  \code{NOT_INGESTED}, \code{NO_FILE}, \code{NO_DOCUMENTS}, \code{UNCHANGED},
  \code{UNKNOWN_KIND}; or a name is taken: \code{EMAIL_TAKEN}, \code{DUPLICATE_NAME},
  \code{SLUG_TAKEN}, \code{ALREADY_ADOPTED}, \code{MESSAGE_ID_TAKEN} \\
413 & \code{FILE_TOO_LARGE}, \code{QUOTA_EXCEEDED}, \code{MESSAGE_TOO_LARGE} \\
429 & \code{RATE_LIMITED} (with a \code{Retry-After} header), \code{BUDGET_EXCEEDED} (the
  caller's daily budget for model calls) \\
500 & \code{SERVER_ERROR}; the traceback goes to the log, never to the client. A value that
  passed the API's own validation and breaks a constraint of the database --- other than a
  taken name, which is a 409 --- ends here too, on purpose: it shows a gap in that validation,
  which a 400 would hide \\
502 & \code{MODEL_ERROR} (a prompt test run whose model call failed) \\
503 & \code{MODELS_UNAVAILABLE} (the gateway's list of models cannot be read) \\
\bottomrule
\end{tabularx}
\end{table}

The assistant service wraps its errors differently, because the chat runtime in the console
expects that shape (\code{errorMiddleware} in the assistant service's
\file{middleware/errors.js}):

\begin{codeblock}
{"message": "Too many questions in a row — wait a moment",
 "error": {"code": "RATE_LIMITED", "message": "Too many questions in a row — wait a moment", "retryAfterS": 41}}
\end{codeblock}

\code{retryAfterS} appears only on a 429 and is repeated in a \code{Retry-After} header. The
codes of the assistant service are listed in \cref{sec:api:assistant}.

\subsection{Lists: paging and filters}\label{sec:api:paging}

A list route answers an object with the rows under a plural name (\code{companies},
\code{reports}) and, where it pages, the count of all matching rows as \code{total}. Only the
long lists page; the others return everything up to a fixed cap (\cref{tab:api:paging}). Two
readers of numbers are in use, and they treat a bad value differently: \code{query_int} refuses
a value outside its range with 400 \code{INVALID_QUERY}, while the jobs, audit and units routes
clamp it into the range and refuse only what is not a number.

\begin{table}[H]
\caption{How the list routes page}\label{tab:api:paging}
\begin{tabularx}{\textwidth}{@{}L{55mm}L{27mm}Y@{}}
\toprule
\hd{List} & \hd{\code{limit}: default, most} & \hd{Paging} \\
\midrule
\code{/api/documents} & 500, 500 & \code{limit}, \code{offset}; answers \code{total}; out of range: 400 \\
\code{/api/obligations} & 500, 2000 & \code{limit}, \code{offset}; answers \code{total}; out of range: 400 \\
\code{/api/usage/calls} & 50, 200 & \code{limit}, \code{offset}; answers \code{total}; out of range: 400 \\
\code{/api/corpus/documents/{id}/units} & 200, 1000 & \code{limit}, \code{offset}, clamped; answers \code{total}, \code{offset}, \code{limit} \\
\code{/api/jobs} & 100, 500 & \code{limit} only, clamped; newest first \\
\code{/api/audit} & 100, 500 & a cursor: \code{before} and \code{beforeId}, taken from the \code{next} of the page before \\
\code{/api/documents/{id}/text} & 20\,000, 100\,000 & a window of characters: \code{offset}, \code{limit}; answers \code{total} \\
companies, chat threads, reports, work products, actions & --- & none; at most 500, 300, 200, 300 and 500 rows \\
\bottomrule
\end{tabularx}
\end{table}

Filters are query parameters named after the field they filter (\code{companyId},
\code{status}, \code{kind}, \code{jurisdiction}). \code{q} is a text search wherever it
appears. \code{scope=own} narrows an administrator's list of companies or documents to the rows
that are the administrator's own, which is what the chat assistant asks for; \code{scope=all}
and \code{userId} on the usage routes widen an administrator's view the other way. A filter
value outside its list is ignored by the corpus and obligation lists and refused with 400
\code{INVALID_QUERY} by the jobs and usage lists; in the list of reports an unknown
\code{status} matches nothing.

\subsection{Languages}\label{sec:api:languages}

The API itself has one language: its error sentences are English, and a client that shows its
own wording keys it on \code{code}. Where a route produces text for a reader, the request says
in which of the product's two languages. A report and a work product take \code{language},
\code{lt} or \code{en}, Lithuanian when absent, and keep it for life. The export of the action
register takes \code{language} as a query parameter. Obligation extraction takes
\code{language: "en"} and is Lithuanian otherwise. A chat turn takes \code{locale} in its body;
without it the \code{Accept-Language} header decides, English when it begins with \code{en} and
Lithuanian in every other case (\code{turnLanguage} in the assistant service's
\file{routes/chat.js}). The \code{language} of an act in the corpus is a fact about its text,
not a choice of output.

\subsection{Limits}\label{sec:api:limits}

\paragraph{Rate limits.} Two things are rate-limited, both with a fixed window per key. The
login counts attempts per client address --- \code{LOGIN_RATE_PER_IP}, 60 by default, in
\code{LOGIN_WINDOW_S}~=~60 seconds --- and per account: \code{EMAIL_LIMIT}~=~10 attempts in
\code{EMAIL_WINDOW_S}~=~600 seconds, a window the password check of
\route{POST}{/api/auth/change-password} shares. The windows are rows of the table
\code{rate_limit_windows} (\file{common/ratelimit.py}), so every process of the API counts
together, and a refused attempt still counts. The chat allows \code{RATE_LIMIT_TURNS}~=~12
turns per user in \code{RATE_LIMIT_WINDOW_MS}~=~60 seconds, counted in the memory of the
assistant service; there a refused turn does not count. Both answer 429 \code{RATE_LIMITED}
with a \code{Retry-After} header. No other route is rate-limited.

\paragraph{The daily budget.} What bounds the costly routes is money, not request counts. Every
user has a daily budget for model calls in US dollars (\code{LLM_DAILY_BUDGET_USD}, \$10.00 by
default, or the account's own \code{dailyBudgetUsd}; an administrator's is unlimited unless
set). A chat turn and a prompt test run are refused with 429 \code{BUDGET_EXCEEDED} once
nothing remains; starting a report, approving its scope and starting a work product need
\code{MIN_BUDGET_TO_START_USD}~=~\$0.25 left. The citation checker is the one route that
degrades instead of refusing. \Cref{ch:models} describes how calls are priced and counted.

\paragraph{Sizes.} The ingress caps no body; each service applies its own limits. The API
reads a body other than an uploaded file up to \code{DATA_UPLOAD_MAX_MEMORY_SIZE}~=~8~MiB
(\file{settings.py}); Django refuses a larger one before any route sees it. A company
document may be 50~MiB (\code{MAX_FILE_BYTES}) and a file for the corpus 60~MiB
(\code{MAX_UPLOAD_BYTES}); beyond that the upload answers 413 \code{FILE_TOO_LARGE}. An account
may store \code{UPLOAD_QUOTA_FILES}~=~1000 documents and \code{UPLOAD_QUOTA_BYTES}~=~2~GiB in
all; beyond that, 413 \code{QUOTA_EXCEEDED}. A stored chat message may be 400\,000 characters
as JSON (413 \code{MESSAGE_TOO_LARGE}). The assistant service accepts 512~kB of JSON for a chat
turn (413 \code{PAYLOAD_TOO_LARGE}). Text fields have their own caps, given in the tables where
they matter.

\begin{codebox}{Strict reading of JSON}
The one reader of request bodies, \code{get_json_object} in \file{common/http.py}, accepts a
JSON \emph{object} and nothing else. It refuses, as if no body had been sent:
\begin{itemize}
  \item an array or a bare value at the top;
  \item \code{NaN} and \code{Infinity}, spelled out or reached by overflow (\code{1e999});
  \item a NUL character or a lone UTF-16 surrogate in any key or string --- no database text
    can hold them;
  \item nesting deeper than \code{MAX_JSON_DEPTH}~=~32.
\end{itemize}
Most routes then answer 400 \code{INVALID_BODY}; a route whose body is optional reads a refused
body as an empty one. Fields are read with typed readers: \code{strict_bool} takes a JSON
boolean and never the string \code{"false"}, \code{bounded_int} takes a whole number within its
range and never a string. A key a route does not know is ignored. In front of all routes,
\code{RequestGuardMiddleware} (\file{common/middleware.py}) answers 400 \code{BAD_REQUEST} to a
NUL character in the path or the query, and turns a value the database cannot take, one out
of range or malformed, into 400
\code{INVALID_VALUE} instead of a 500. A constraint of the database that a value breaks is
left a 500 (\cref{tab:api:codes}).
\end{codebox}


% ============================================================
\section{The routes of the API}\label{sec:api:routes}
% ============================================================

\file{urls.py} holds 83 URL patterns: 77 under \code{/api/}, listed in this section, and 6
under \code{/internal/assistant/}, listed in \cref{sec:api:internal}. Counted by method, the 83
patterns serve 119 method--path pairs. Each table below has one entry per pattern: the methods
and the path on the first line, then who may call (\cref{tab:api:who}) and what the route does.
In a path, \code{{id}} stands for the id of the row. \enquote{Errors} names the refusals
particular to the route.

Two boxes first sort the routes by the question this document asks of every part of the
system: who is acting.

\begin{aibox}{Routes that call a model before they answer}
Almost every route that needs a model only queues a job and answers at once; the model is then
called by the worker (\cref{ch:arch}). Five kinds of request wait for a model themselves:
\begin{itemize}
  \item \route{POST}{/api/corpus/check-citations} --- one call of the judge behind the citation
    check, when the request says what a provision is claimed to be about (\cref{ch:company});
  \item \route{POST}{/api/prompts/{slug}/test} --- one call with a prompt under test
    (\cref{ch:prompts});
  \item \route{POST}{/api/corpus/search} and \route{GET}{/api/documents/search} --- one
    embedding of the query (\cref{ch:corpus});
  \item \route{POST}{/api/obligations} and a \code{PATCH} that changes an obligation's title or
    text --- one embedding of the obligation (\cref{ch:obligations});
  \item \route{POST}{/api/assistant/chat} --- the chat turn itself (\cref{ch:chat}).
\end{itemize}
\end{aibox}

\begin{humanbox}{Routes through which a person decides}
The decisions the system never takes itself each enter through one route, are stamped with who
decided and when, and are written to the audit trail: approving the scope of a report
(\route{POST}{/api/reports/{id}/scope} with \code{approved: true}); overriding a finding and
writing its cause (\route{PATCH}{/api/reports/{id}/findings/{fid}}); signing a report off
(\route{PATCH}{/api/reports/{id}}); approving, correcting or rejecting an obligation
(\route{PATCH}{/api/obligations/{id}}, \route{POST}{/api/obligations/review}); approving a work
product (\route{PATCH}{/api/programme/{id}}); marking a credential as verified
(\route{PATCH}{/api/companies/{id}/credentials/{cid}}, or the \code{POST} that files it with
that status); closing an action by hand, with a note
(\route{PATCH}{/api/actions/{id}}).
\end{humanbox}

\subsection{Health}\label{sec:api:health}

One public probe says whether the API can do its work (\file{ops/api/views.py}).

\begin{xltabular}{\textwidth}{@{}L{19mm}Y@{}}
\caption{Routes: health}\label{tab:api:health}\\
\toprule \hd{Who} & \hd{What the route does} \\ \midrule \endfirsthead
\toprule \hd{Who} & \hd{What the route does} \\ \midrule \endhead
\bottomrule \endlastfoot
\apiroute{GET HEAD}{/api/health}
public & Readiness: the database answers a trivial query and the upload directory is writable.
  200 with \code{status: "ok"}, \code{service: "legalops-django"}, \code{database},
  \code{uploads}; or 503 with \code{status: "error"} and \code{reason}, the first check that
  failed (\code{database} or \code{uploads}) --- which check, never why: the cause is in the
  log. \\
\end{xltabular}

\subsection{Authentication}\label{sec:api:authroutes}

The caller's own session (\file{users/api/auth_views.py}). A \code{user} in an answer is the
account as \code{serialize_user} writes it: \code{id}, \code{email}, \code{displayName},
\code{role}, \code{active}, \code{dailyBudgetUsd}, \code{createdAt}, \code{lastLoginAt}.

\begin{xltabular}{\textwidth}{@{}L{19mm}Y@{}}
\caption{Routes: authentication}\label{tab:api:auth}\\
\toprule \hd{Who} & \hd{What the route does} \\ \midrule \endfirsthead
\toprule \hd{Who} & \hd{What the route does} \\ \midrule \endhead
\bottomrule \endlastfoot
\apiroute{POST}{/api/auth/login}
public & Body \code{email}, \code{password}. Answers \code{token} and \code{user}, and sets the
  gate cookie. The e-mail is matched without regard to case. An unknown e-mail, a wrong
  password and a deactivated account get the same answer, 401 \code{INVALID_CREDENTIALS}, at
  the same cost in time. Successes and failures are audited. \emph{Errors:} 400
  \code{MISSING_FIELDS}; 429 \code{RATE_LIMITED}. \\
\addlinespace
\apiroute{GET}{/api/auth/me}
signed in & The caller's own \code{user}. The console calls it at every start; the assistant
  service calls it to turn a token into a user. \\
\addlinespace
\apiroute{GET}{/api/auth/checkauth/admin}
admin & The gate the ingress asks before a request to \code{/dbgate/}. The session is taken
  from the token or, because a tool opened in the browser sends only cookies, from the gate
  cookie. 200 with the \code{user} for an active administrator; 401 \code{AUTH_REQUIRED}; 403
  \code{FORBIDDEN} for any other role. \\
\addlinespace
\apiroute{POST}{/api/auth/logout}
signed in & Ends the session of the presented token and expires the gate cookie. Other devices
  stay signed in. \\
\addlinespace
\apiroute{POST}{/api/auth/logout-all}
signed in & Ends every session of the caller, the presented one included. Audited. \\
\addlinespace
\apiroute{POST}{/api/auth/change-password}
signed in & Body \code{currentPassword}, \code{newPassword}. Every \emph{other} session of the
  account ends. A password has at least 10 characters and at most 72 bytes. \emph{Errors:} 400
  \code{WRONG_PASSWORD} (audited), \code{WEAK_PASSWORD}; 429 \code{RATE_LIMITED}. \\
\end{xltabular}

\subsection{Users}\label{sec:api:users}

The directory of accounts, for administrators (\file{users/api/users_views.py}). Every change
is audited. Two guards hold on every change of role or state and on deletion: the last active administrator can be neither
demoted, deactivated nor deleted (400 \code{LAST_ADMIN}), and nobody can do any of these to
their own account (400 \code{SELF_LOCKOUT}).

\begin{xltabular}{\textwidth}{@{}L{19mm}Y@{}}
\caption{Routes: users}\label{tab:api:users}\\
\toprule \hd{Who} & \hd{What the route does} \\ \midrule \endfirsthead
\toprule \hd{Who} & \hd{What the route does} \\ \midrule \endhead
\bottomrule \endlastfoot
\apiroute{GET POST}{/api/users}
admin & \code{GET}: every account, deactivated ones included, oldest first, as \code{users}.
  \code{POST}: body \code{email}, \code{displayName}, \code{password} and \code{role}
  (\code{user} when absent); answers 201 with the account. The e-mail is stored in lower case.
  \emph{Errors:} 400 \code{INVALID_EMAIL}, \code{MISSING_FIELDS}, \code{INVALID_ROLE},
  \code{WEAK_PASSWORD}; 409 \code{EMAIL_TAKEN} (in any casing). \\
\addlinespace
\apiroute{GET PATCH DELETE}{/api/users/{id}}
admin & \code{PATCH}: any of \code{displayName}, \code{role}, \code{active},
  \code{dailyBudgetUsd} (0 to 10\,000; \code{null} means the default and 0 means unlimited).
  \code{DELETE} deactivates the account and ends its sessions; the rows stay. \code{DELETE}
  with \code{purge=true} deletes the account for good together with its companies, documents
  and conversations, and answers what it removed. \emph{Errors:} 400 \code{EMPTY_UPDATE},
  \code{INVALID_FIELD}; 409 \code{STILL_RUNNING} (a purge while a report of one of the
  account's companies is queued or running). \\
\addlinespace
\apiroute{POST}{/api/users/{id}/reset-password}
admin & Body \code{password}. Sets a new password for the account and ends all its
  sessions. \emph{Errors:} 400 \code{WEAK_PASSWORD}. \\
\end{xltabular}

\subsection{Companies}\label{sec:api:companies}

The company profile and what hangs on it: credentials, adopted standards and roles
(\file{companies/api/company_views.py}, \file{companies/api/standard_views.py},
\file{companies/api/role_views.py}). What
these fields mean to a report is the subject of \cref{ch:company}. Every write is audited.

\begin{xltabular}{\textwidth}{@{}L{19mm}Y@{}}
\caption{Routes: companies, credentials, standards and roles}\label{tab:api:companies}\\
\toprule \hd{Who} & \hd{What the route does} \\ \midrule \endfirsthead
\toprule \hd{Who} & \hd{What the route does} \\ \midrule \endhead
\bottomrule \endlastfoot
\apiroute{GET POST}{/api/companies}
owner & \code{GET}: the caller's companies as \code{companies}, most recently updated first,
  at most 500. \code{q} matches word by word against the name and the registration code;
  punctuation, legal forms (UAB, LTD) and the word \enquote{company} are skipped.
  \code{scope=own} narrows an administrator to their own. \code{POST}: \code{name} and any
  profile fields; answers 201 with the company; the caller becomes its owner. \emph{Errors:}
  400 \code{MISSING_FIELDS}, \code{INVALID_FIELD}. \\
\addlinespace
\apiroute{GET PUT DELETE}{/api/companies/{id}}
owner & \code{GET}: the profile with its \code{credentials}, \code{standards}, \code{roles}
  and \code{documents}. \code{PUT} is a partial update: any subset of the profile fields.
  \code{DELETE} removes the company with its credentials, adopted standards, roles, reports and
  work products; its documents are detached, not deleted. \emph{Errors:} 400
  \code{INVALID_FIELD}; 409 \code{STILL_RUNNING} (a report or a work product is in work). \\
\addlinespace
\apiroute{GET POST}{/api/companies/{id}/credentials}
owner & A credential is a licence, certificate, permit or registration the company holds.
  \code{POST}: \code{kind} (\code{certificate}, \code{license}, \code{permit},
  \code{registration}), \code{name}, and optionally \code{issuer}, \code{identifier},
  \code{validFrom}, \code{validTo}, \code{scope}, \code{notes}, \code{status},
  \code{evidenceDocumentId}. A new credential is \code{claimed} unless the body says
  \code{verified}. \emph{Errors:} 400 \code{MISSING_FIELDS}, \code{INVALID_FIELD}. \\
\addlinespace
\apiroute{PATCH DELETE}{/api/companies/{id}/credentials/{cid}}
owner & \code{PATCH}: any subset of the same fields. Setting \code{status} to \code{verified}
  stamps who verified it and when; \code{claimed} clears the stamp. The evidence must be one of
  the company's own documents. \\
\addlinespace
\apiroute{GET POST}{/api/companies/{id}/standards}
owner & The standards the company has adopted. \code{POST}: \code{documentId} (a standard of
  the corpus), \code{status} (\code{certified}, \code{implementing}, \code{planned};
  \code{implementing} when absent), \code{scope}, \code{certificateId} (one of the company's
  credentials), \code{notes}. \emph{Errors:} 400 \code{MISSING_FIELDS}, \code{INVALID_FIELD}; 409
  \code{ALREADY_ADOPTED}. \\
\addlinespace
\apiroute{PATCH DELETE}{/api/companies/{id}/standards/{sid}}
owner & \code{PATCH}: any of \code{status}, \code{scope}, \code{certificateId} (\code{null}
  unlinks it), \code{notes}. The standard itself never changes through these routes. \\
\addlinespace
\apiroute{GET POST}{/api/companies/{id}/roles}
owner & A role is a position in the company, such as the compliance officer, described for the
  chat assistant. \code{POST}: \code{name}, \code{context} (at most 6\,000 characters),
  \code{sources} (at most 50 entries of \code{documentId}, \code{title}, \code{section}).
  A source that names one of the company's documents, by id or by title, is stored and
  returned with both.
  \emph{Errors:} 400 \code{MISSING_FIELDS}, \code{INVALID_FIELD}; 409 \code{DUPLICATE_NAME}
  (within the company). \\
\addlinespace
\apiroute{GET PATCH DELETE}{/api/companies/{id}/roles/{rid}}
owner & \code{PATCH}: any of \code{name}, \code{context}, \code{sources}. Deleting a role
  unlinks it from the chat threads that adopted it. \\
\end{xltabular}

The profile fields a company takes are \code{name}, \code{registrationCode}, \code{country},
\code{legalForm}, \code{naceCodes}, \code{sector}, \code{sizeCategory}, \code{employees},
\code{annualTurnoverEur}, \code{activities}, \code{usesAi}, \code{providesAi},
\code{aiDescription}, \code{processesPersonalData}, \code{isPublicSector},
\code{institutionType}, \code{financialMarket} and \code{attributes}. The four yes-or-no fields
take \code{true}, \code{false} or \code{null}, and \code{null} means \emph{unknown}: the API
never guesses one, and refuses the string \code{"false"} rather than read it as true.

\subsection{Documents}\label{sec:api:documents}

A company's documents: upload, search and reading (\file{companies/api/document_views.py};
\cref{ch:company}). A document belongs to the user who uploaded it and may be attached to a
company at upload or later.

\begin{xltabular}{\textwidth}{@{}L{19mm}Y@{}}
\caption{Routes: documents}\label{tab:api:documents}\\
\toprule \hd{Who} & \hd{What the route does} \\ \midrule \endfirsthead
\toprule \hd{Who} & \hd{What the route does} \\ \midrule \endhead
\bottomrule \endlastfoot
\apiroute{GET POST}{/api/documents}
owner & \code{GET}: \code{documents} and \code{total}, newest first, without their text;
  filters \code{companyId}, \code{q} (the title), \code{scope=own}. \code{POST} is
  \code{multipart/form-data}: \code{file}, and optionally \code{companyId}, \code{title} (the
  file name when absent), \code{kind} (\code{policy}, \code{procedure}, \code{contract},
  \code{register}, \code{other}). The text is extracted before the answer, 201 with the
  document; its \code{extractionStatus} is \code{done}, \code{empty} (no text, as in a scanned
  PDF) or \code{failed}. When there is text, the job \code{documents.index} is queued.
  \emph{Errors:} 400 \code{MISSING_FILE}, \code{INVALID_FIELD}; 413 \code{FILE_TOO_LARGE},
  \code{QUOTA_EXCEEDED}. \\
\addlinespace
\apiroute{GET}{/api/documents/search}
owner & Passages of the caller's documents that match \code{q} (required, at most 1\,000
  characters), as \code{hits} with the document, the heading, the text and \code{charStart}.
  \code{k} is the number of hits, 1 to 30, 8 when absent; \code{companyId}, \code{documentIds}
  (comma-separated, at most 50) and \code{scope=own} narrow the search. \emph{Errors:} 400
  \code{INVALID_QUERY}. \\
\addlinespace
\apiroute{GET PATCH DELETE}{/api/documents/{id}}
owner & \code{GET}: the metadata with the \code{outline} (headings and their offsets) and the
  \code{inventory}. \code{PATCH}: \code{companyId} (\code{null} detaches the document),
  \code{title}, \code{kind}. \code{DELETE} removes the row and then the file. \emph{Errors:}
  400 \code{EMPTY_UPDATE}, \code{INVALID_FIELD}. \\
\addlinespace
\apiroute{GET}{/api/documents/{id}/text}
owner & A window of the extracted text: \code{offset} (0 when absent) and \code{limit}
  (20\,000 when absent, at most 100\,000). Answers \code{text}, \code{offset}, \code{limit},
  \code{total} and \code{extractionStatus}. \emph{Errors:} 400 \code{INVALID_QUERY}. \\
\end{xltabular}

\subsection{Chat}\label{sec:api:chat}

The API keeps the chat's threads and messages (\file{chat/api/views.py}); the turn itself is
answered by the assistant service (\cref{sec:api:assistant}, \cref{ch:chat}). Threads are
visible to their own user only.

\begin{xltabular}{\textwidth}{@{}L{19mm}Y@{}}
\caption{Routes: chat threads and messages}\label{tab:api:chat}\\
\toprule \hd{Who} & \hd{What the route does} \\ \midrule \endfirsthead
\toprule \hd{Who} & \hd{What the route does} \\ \midrule \endhead
\bottomrule \endlastfoot
\apiroute{GET}{/api/chat/models}
signed in & The models a conversation may answer with: \code{models}, the gateway's chat
  models with \code{id}, \code{name} and \code{pricing}, and \code{default}, the model of the
  active version of \prompt{chat.assistant}. The list is empty while the gateway cannot be
  reached. \\
\addlinespace
\apiroute{GET POST}{/api/chat/threads}
owner only & \code{GET}: the caller's threads that are not archived, most recently updated
  first, at most 300; filter \code{companyId}. \code{POST}: optional \code{companyId} and
  \code{title}; answers 201 with the thread. \\
\addlinespace
\apiroute{GET PATCH DELETE}{/api/chat/threads/{id}}
owner only & \code{PATCH}: \code{title} (at most 200 characters), \code{companyId}
  (\code{null} unlinks the company; a change drops the role), \code{roleId} (a role of the
  thread's company), \code{model} (one of the ids of \code{/api/chat/models}, or \code{null}
  for the default). \code{DELETE} archives the thread; the rows stay. \emph{Errors:} 400
  \code{UNKNOWN_MODEL}, \code{EMPTY_UPDATE}; 503 \code{MODELS_UNAVAILABLE}. \\
\addlinespace
\apiroute{GET POST}{/api/chat/threads/{id}/messages}
owner only & \code{GET}: the history, a bare list of rows \code{id}, \code{parent_id},
  \code{format}, \code{content}. \code{POST} stores or replaces one message by its id within
  the thread and answers \code{ok}, \code{created} and \code{autoTitle} (whether the thread
  still carries the title it gave itself). Only the user's own messages may be written here;
  the assistant's replies come through the internal door. \emph{Errors:} 403
  \code{ROLE_NOT_ALLOWED}; 409 \code{MESSAGE_ID_TAKEN}; 413 \code{MESSAGE_TOO_LARGE}. \\
\end{xltabular}

\subsection{The legal corpus}\label{sec:api:corpus}

Acts, their versions and provisions, search, and the citation checker
(\file{corpus/api/views.py}; \cref{ch:corpus}). Reading is open to everyone signed in; a user
gets 404 for an act that is not enabled. Writing is for administrators and audited. No route
here ingests anything itself: the routes that process an act queue a job and answer 202 with
it (\cref{sec:api:client}).

\begin{xltabular}{\textwidth}{@{}L{19mm}Y@{}}
\caption{Routes: the legal corpus}\label{tab:api:corpus}\\
\toprule \hd{Who} & \hd{What the route does} \\ \midrule \endfirsthead
\toprule \hd{Who} & \hd{What the route does} \\ \midrule \endhead
\bottomrule \endlastfoot
\apiroute{GET POST}{/api/corpus/documents}
signed in; \code{POST}: admin & \code{GET}: the catalogue as \code{documents}, each act with its
  current version, status, last error and number of obligations; filters \code{q},
  \code{jurisdiction}, \code{docType}, \code{domain}. \code{POST} adds an act: \code{title},
  \code{jurisdiction} (\code{EU}, \code{LT}, \code{INTL}), \code{sourceKind} (\code{cellar},
  \code{etar}, \code{eseimas}, \code{lb}, \code{url}, \code{upload}; \code{upload} when absent)
  and, as the source needs, \code{externalId} or an https \code{sourceUrl}; also \code{shortTitle}, \code{titleEn},
  \code{docType}, \code{language}, \code{domains}, \code{appliesTo},
  \code{applicabilityNotes}, \code{status}, \code{enabled}. An act with a source that can be
  fetched queues its ingestion at once; an \code{upload} act waits for its file.
  \emph{Errors:} 400 \code{INVALID_FIELD}; 409 \code{SLUG_TAKEN}. \\
\addlinespace
\apiroute{GET PATCH DELETE}{/api/corpus/documents/{id}}
signed in; \code{PATCH}, \code{DELETE}: admin & \code{GET}: the act with its \code{versions}
  and \code{latestJob}, and what the models read from the current version's own text:
  \code{scopeReading} (\code{types}, \code{reasoning}, \code{readAt}) and
  \code{sanctionsReading} (\code{tiers} --- the number of tiers, \code{note}, \code{readAt}), each
  \code{null} until it is read. \code{PATCH} edits the same fields. \code{DELETE} removes the act with
  every version, unit, chunk and stored file. \\
\addlinespace
\apiroute{POST}{/api/corpus/documents/{id}/refresh}
admin & Queues a fetch and ingestion. \code{force} processes the text again even when the
  source is unchanged; \code{structure} (\code{auto}, \code{rules}, \code{ai}, \code{pages};
  for a standard only) says how the standard's structure is read this time. Answers 202 with
  the job --- the one already queued or running, when there is one. \emph{Errors:} 400
  \code{INVALID_FIELD}; 409 \code{NO_FILE} (an \code{upload} act without a file). \\
\addlinespace
\apiroute{POST}{/api/corpus/documents/{id}/upload}
admin & \code{multipart/form-data} with \code{file}: the file becomes a new version and its
  ingestion is queued; 202 with the job. \emph{Errors:} 400 \code{MISSING_FILE}; 409
  \code{UNCHANGED} (the same bytes as the current version); 413 \code{FILE_TOO_LARGE}. \\
\addlinespace
\apiroute{GET}{/api/corpus/documents/{id}/units}
signed in & The provisions of the current version, or of \code{version}, in the order of the
  act, with shortened texts: \code{units}, \code{total}, \code{offset}, \code{limit}. Filters
  \code{type} (a comma-separated list of unit types) and \code{q}. \\
\addlinespace
\apiroute{GET}{/api/corpus/units/{id}}
signed in & One provision in full, with its act, \code{versionNo}, \code{isCurrent} and
  \code{citationUrl}, the link to the official text. \\
\addlinespace
\apiroute{GET}{/api/corpus/provision}
signed in & The exact text of one provision, found by reference: \code{document} (a short
  title, slug, id or official number), \code{article} (or \code{clause}, for a standard),
  \code{paragraph}, \code{type}. An inserted article is found however its number is written
  (\enquote{12-1}, \enquote{12¹}, \enquote{12(1)}). 404 when the corpus holds no such
  provision. \\
\addlinespace
\apiroute{POST}{/api/corpus/search}
signed in & Search over the current versions. Body \code{query}, and optionally \code{k} (1 to
  30, 8 when absent), \code{documents} (at most 50 acts, each named as any lookup of an act
  takes it --- id, slug, short title or CELEX number, whatever the case; a name that means no
  act narrows the search to nothing), \code{jurisdiction}, \code{kind} (\code{acts} or
  \code{standards}). Answers \code{hits}.
  \emph{Errors:} 400 \code{MISSING_FIELDS}, \code{INVALID_FIELD}. \\
\addlinespace
\apiroute{POST}{/api/corpus/check-citations}
signed in & The citation checker. Body \code{citations}, 1 to 50 items of \code{act},
  \code{article} and optionally \code{paragraph}, \code{point}, \code{claimedHeading}, each at
  most 300 characters. Answers \code{results}, one per citation in order: \code{status}
  (\code{found}, \code{not_found}, \code{act_not_in_corpus}), the provision found,
  \code{headingMatch} (\code{match}, \code{mismatch}, \code{unknown}) and \code{judgedBy}
  (\code{model}, \code{overlap} or \code{null}). When the caller's budget is spent, or the
  model call fails, the claim is judged by word overlap instead of the model; the check itself
  never stops. \emph{Errors:}
  400 \code{INVALID_FIELD}. \\
\addlinespace
\apiroute{POST}{/api/corpus/reembed}
admin & Queues the re-embedding of one act (\code{documentId}) or, without it, of all; 202
  with the job. Needed only after the embedding model changes. \\
\addlinespace
\apiroute{POST}{/api/corpus/refresh-all}
admin & Queues a check of every source for a new version; 202 with the job. \\
\end{xltabular}

\subsection{Obligations}\label{sec:api:obligations}

The register of obligations and its review (\file{corpus/api/obligation_views.py};
\cref{ch:obligations}). Everyone signed in may read; only an administrator changes an
obligation, and every change is audited.

\begin{xltabular}{\textwidth}{@{}L{19mm}Y@{}}
\caption{Routes: obligations}\label{tab:api:obligations}\\
\toprule \hd{Who} & \hd{What the route does} \\ \midrule \endfirsthead
\toprule \hd{Who} & \hd{What the route does} \\ \midrule \endhead
\bottomrule \endlastfoot
\apiroute{GET POST}{/api/obligations}
signed in; \code{POST}: admin & \code{GET}: \code{obligations} and \code{total}. Filters
  \code{document}, \code{status} (\code{proposed}, \code{approved}, \code{rejected}; one, or
  several with commas between), \code{category}, \code{q}, and \code{verification}, the judge's verdict: \code{supported},
  \code{partly}, \code{unsupported}, \code{doubt} (the last two together) or \code{unchecked}.
  \code{order=review} puts the doubted first and the gravest first among them. \code{POST}
  enters an obligation by hand: \code{title}, \code{text}, optionally \code{documentId} and the
  editable fields below; it is stored with the source \code{standard}, approved at once.
  \emph{Errors:} 400 \code{INVALID_QUERY}, \code{MISSING_FIELDS}, \code{INVALID_FIELD}. \\
\addlinespace
\apiroute{POST}{/api/obligations/review}
admin & The bulk decision: \code{ids} (1 to 500) and \code{status}. Each obligation is stamped
  with the reviewer; one audit row lists them. Answers \code{updated} and \code{status}.
  \emph{Errors:} 400 \code{INVALID_FIELD}. \\
\addlinespace
\apiroute{GET PATCH DELETE}{/api/obligations/{id}}
signed in; \code{PATCH}, \code{DELETE}: admin & \code{PATCH}: any of \code{title},
  \code{titleEn}, \code{text}, \code{textEn}, \code{addressee}, \code{conditions},
  \code{evidence}, \code{deadline}, \code{notes}, \code{category}, \code{severity},
  \code{appliesTo}, \code{status}. Setting \code{status} stamps the reviewer. A changed title
  or text is embedded again and loses the judge's verdict, which was given on other words.
  \emph{Errors:} 400 \code{INVALID_FIELD}. \\
\addlinespace
\apiroute{POST}{/api/corpus/documents/{id}/extract-obligations}
admin & Queues the extraction of obligations from the act's current version; \code{language}.
  202 with the job, or with the one already at work. \emph{Errors:} 409 \code{NOT_INGESTED}
  (the act has no version yet). \\
\addlinespace
\apiroute{POST}{/api/corpus/documents/{id}/verify-obligations}
admin & Queues the judge over the obligations of the current version that it has not judged
  yet --- all of them with \code{again: true}. 202 with the job, or with the one already at
  work. \emph{Errors:} 400 \code{INVALID_FIELD}; 409 \code{NOT_INGESTED}. \\
\end{xltabular}

\subsection{Reports}\label{sec:api:reports}

Starting a compliance report, the scope review, the review of findings, the sign-off and the
exports (\file{analysis/api/views.py}; \cref{ch:report}). A report follows its company: its
owner and administrators see it. One report per company is in work at a time.

\begin{xltabular}{\textwidth}{@{}L{19mm}Y@{}}
\caption{Routes: reports}\label{tab:api:reports}\\
\toprule \hd{Who} & \hd{What the route does} \\ \midrule \endfirsthead
\toprule \hd{Who} & \hd{What the route does} \\ \midrule \endhead
\bottomrule \endlastfoot
\apiroute{GET POST}{/api/reports}
owner & \code{GET}: \code{reports}, newest first, at most 200, each with its job's
  \code{jobStatus} and \code{progress}; filters \code{companyId} and \code{status} (the
  report's state), and \code{scope=own}, which narrows an administrator to the reports of
  their own companies --- all applied before the list is cut to its 200. \code{POST} starts a report:
  \code{companyId}, \code{language}, and optionally \code{documentIds} (at most 200 enabled
  acts, to restrict the report to them), \code{review} (stop for the scope review) and
  \code{internal} (also assess the company's own requirements). Answers 201 with the queued
  report; the job is \code{report.analyze}. \emph{Errors:} 400 \code{INVALID_FIELD}; 409
  \code{CORPUS_EMPTY}, \code{ALREADY_RUNNING} (also while a report of the company waits for
  its scope); 429 \code{BUDGET_EXCEEDED}. \\
\addlinespace
\apiroute{GET}{/api/reports/estimate}
owner & What a report of \code{companyId} would cost now, as an upper bound in US dollars,
  with the caller's \code{budget}; \code{documentIds} and \code{internal=true} as for the
  start. \emph{Errors:} 400 \code{INVALID_FIELD} for \code{documentIds} the start would
  refuse (an act that is not in the corpus, or disabled); 404 \code{NOT_FOUND}. \\
\addlinespace
\apiroute{GET PATCH DELETE}{/api/reports/{id}}
owner & \code{GET}: the whole report --- \code{acts}, \code{findings}, \code{actions},
  \code{citations}, \code{summary}, \code{stats}, the approvals, and the job's status and
  progress. \code{PATCH} is the sign-off: \code{approved} (\code{true} or \code{false}) and
  \code{note}; the answer adds \code{register}, what the sign-off moved in the action
  register. \code{DELETE} first withdraws a sign-off. \emph{Errors:} 400
  \code{MISSING_FIELDS}; 409 \code{NOT_FINISHED} (sign-off), \code{STILL_RUNNING} (delete). \\
\addlinespace
\apiroute{POST}{/api/reports/{id}/cancel}
owner & Stops a report that is queued or running, and cancels at once one that waits for its
  scope. \emph{Errors:} 409 \code{NOT_CANCELLABLE}. \\
\addlinespace
\apiroute{POST}{/api/reports/{id}/scope}
owner & The scope review of a report that waits in \code{scope_review}. Body \code{acts}, a
  map from an act's id in the report to \code{assess} or \code{skip} (an act not named keeps
  what the run decided), and \code{approved}. Without \code{approved} nothing is stored and the
  answer prices the selection: \code{estimate} and \code{budget}. With \code{approved: true}
  the choices are applied, the approval is stamped and audited with every change, and the job
  is queued for the second part; the answer is the report. \emph{Errors:} 400
  \code{INVALID_FIELD}, \code{NOTHING_TO_ASSESS}; 409 \code{NOT_WAITING}; 429
  \code{BUDGET_EXCEEDED}. \\
\addlinespace
\apiroute{PATCH}{/api/reports/{id}/findings/{fid}}
owner & A person's verdict on one finding: \code{overrideStatus} (\code{met}, \code{partial},
  \code{missing}, \code{unknown}, \code{not_applicable}, or \code{null} to clear it),
  \code{reviewerNote} (4\,000 characters), \code{cause} (2\,000). Any of them withdraws a
  sign-off; the report's counts are made again. Audited. \emph{Errors:} 400
  \code{INVALID_FIELD}. \\
\addlinespace
\apiroute{GET}{/api/reports/{id}/export}
owner & The finished report as a file: \code{format} is \code{pdf} (when absent) or
  \code{xlsx}. Audited. \emph{Errors:} 400 \code{INVALID_QUERY}; 409 \code{NOT_FINISHED}. \\
\end{xltabular}

\subsection{The action register and changes in the law}\label{sec:api:actions}

What lives on after a report is signed off (\file{analysis/api/views.py},
\file{analysis/register.py}; \cref{ch:register}). The register itself moves only on a sign-off;
these routes read it and let a person work on single actions.

\begin{xltabular}{\textwidth}{@{}L{19mm}Y@{}}
\caption{Routes: the action register and changes in the law}\label{tab:api:actions}\\
\toprule \hd{Who} & \hd{What the route does} \\ \midrule \endfirsthead
\toprule \hd{Who} & \hd{What the route does} \\ \midrule \endhead
\bottomrule \endlastfoot
\apiroute{GET}{/api/companies/{id}/actions}
owner & The company's action register: \code{actions} (at most 500), the late ones first, then
  the other open ones by due date, then those done and closed; and \code{counts} by state
  (\code{late}, \code{open}, \code{in_progress}, \code{done}, \code{closed}). \\
\addlinespace
\apiroute{GET}{/api/companies/{id}/actions/export}
owner & The register as a workbook (\code{.xlsx}); \code{language}. Audited. \\
\addlinespace
\apiroute{PATCH}{/api/actions/{id}}
owner & A person's hand on one action: \code{status} (\code{open}, \code{in_progress},
  \code{done}, \code{closed}), \code{owner}, \code{dueDate}, \code{note}. Audited.
  \emph{Errors:} 400 \code{INVALID_FIELD}, \code{NOTE_REQUIRED} (closing by hand without a
  note); 409 \code{CLOSED_BY_REPORT} (reopening an action a signed-off report closed). \\
\addlinespace
\apiroute{GET}{/api/companies/{id}/law-changes}
owner & For the company's latest finished report: the assessed acts that have a newer version
  since, their changed provisions, and the findings that rest on them. Answers
  \code{reportId}, \code{reportDate}, \code{acts}. \\
\end{xltabular}

\subsection{Programme}\label{sec:api:programme}

The work products of a company --- kinds \code{policy}, \code{governance} and \code{risks}
(\file{programme/api/views.py}; \cref{ch:programme}). Tenancy follows the company. One product
of a kind is generated at a time per company.

\begin{xltabular}{\textwidth}{@{}L{19mm}Y@{}}
\caption{Routes: work products}\label{tab:api:programme}\\
\toprule \hd{Who} & \hd{What the route does} \\ \midrule \endfirsthead
\toprule \hd{Who} & \hd{What the route does} \\ \midrule \endhead
\bottomrule \endlastfoot
\apiroute{GET POST}{/api/companies/{id}/programme}
owner & \code{GET}: \code{products}, by kind and newest version first, without their content,
  each with its job's status and progress; and \code{kinds}, the three kinds with their titles
  and what each builds on. \code{POST}: \code{kind} and \code{language}; starts the generation
  of a new version (job \code{programme.generate}) and answers 201 with the queued product.
  \emph{Errors:} 400 \code{INVALID_FIELD}; 409 \code{NO_DOCUMENTS} (the company has no
  readable document), \code{ALREADY_RUNNING}; 429 \code{BUDGET_EXCEEDED}. \\
\addlinespace
\apiroute{GET PATCH DELETE}{/api/programme/{id}}
owner & \code{GET}: the product with its \code{content}. \code{PATCH} with
  \code{approved: true} and optionally \code{note} is a person's approval of a draft.
  \code{DELETE} removes a version. \emph{Errors:} 400 \code{MISSING_FIELDS}; 409
  \code{NOT_DRAFT}, \code{RULES_BROKEN} (a blocking rule of the kind is broken; the sentence
  names the rules), \code{STILL_RUNNING} (delete). \\
\addlinespace
\apiroute{PUT}{/api/programme/{id}/content}
owner & A person's edit of a draft, sent whole as \code{content}. The content is checked
  against the rules of its kind without a model call, and the answer is the product as stored,
  with what is derived computed again. \emph{Errors:} 400 \code{MISSING_FIELDS},
  \code{INVALID_FIELD}; 409 \code{NOT_DRAFT}, \code{UNKNOWN_KIND}. \\
\addlinespace
\apiroute{POST}{/api/programme/{id}/revise}
owner & A new draft, one version higher, copied from an approved product, which stays as
  approved; 201. \emph{Errors:} 409 \code{NOT_APPROVED}, \code{NEWER_VERSION}. \\
\addlinespace
\apiroute{POST}{/api/programme/{id}/cancel}
owner & Stops a generation that is queued or running. \emph{Errors:} 409
  \code{NOT_CANCELLABLE}. \\
\addlinespace
\apiroute{GET}{/api/programme/{id}/chart}
owner & The structure of a governance product drawn as an organisation chart,
  \code{image/svg+xml}. 404 when the product has no structure. \\
\addlinespace
\apiroute{GET}{/api/programme/{id}/export}
owner & The product as a file: \code{format} is \code{pdf} for every kind, also \code{docx}
  for a policy and \code{xlsx} for a risk register. A draft exports as well as an approved
  product. Audited. \emph{Errors:} 400 \code{INVALID_QUERY}; 409 \code{NOT_READY}. \\
\end{xltabular}

\subsection{Prompts}\label{sec:api:prompts}

The store of prompts and their versions, for administrators (\file{prompts/api/views.py};
\cref{ch:prompts}). Every write and every test run is audited. The list of models that feeds
the editor's model picker stands in \file{assistant/api/mcp_views.py}.

\begin{xltabular}{\textwidth}{@{}L{19mm}Y@{}}
\caption{Routes: prompts and models}\label{tab:api:prompts}\\
\toprule \hd{Who} & \hd{What the route does} \\ \midrule \endfirsthead
\toprule \hd{Who} & \hd{What the route does} \\ \midrule \endhead
\bottomrule \endlastfoot
\apiroute{GET}{/api/prompts}
admin & Every prompt as \code{prompts}: \code{slug}, \code{name}, \code{kind},
  \code{variables}, \code{activeVersion}, \code{activeModel}, \code{pendingDefault}. \\
\addlinespace
\apiroute{GET}{/api/prompts/{slug}}
admin & One prompt with its \code{versions}, newest first, each with its text. \\
\addlinespace
\apiroute{POST}{/api/prompts/{slug}/versions}
admin & Appends a version: \code{text} (at most 60\,000 characters), \code{model}, and
  optionally \code{temperature} (0 to 2), \code{maxTokens} (1 to 128\,000), \code{note} (500
  characters), \code{activate}. Answers 201 with the version. \emph{Errors:} 400
  \code{MISSING_FIELDS}, \code{INVALID_FIELD}. \\
\addlinespace
\apiroute{POST}{/api/prompts/{slug}/versions/{n}/activate}
admin & Makes version \code{n} the active one; answers the prompt. \\
\addlinespace
\apiroute{POST}{/api/prompts/{slug}/test}
admin & One model call with a stored version (\code{versionNo}; the active one when absent)
  or an unsaved draft (\code{text}, \code{model}) as the instructions and \code{input} (at most
  200\,000 characters) as the data, on the output schema of the prompt's step when the prompt
  has one; \code{variables} fill the
  placeholders. Answers \code{text}, \code{object}, \code{model}, \code{finishReason},
  \code{usage}, \code{costUsd}, \code{durationMs}, \code{renderedInstructions}. A call that answers is
  charged to the caller's budget and logged as a model call; a failed one is only audited.
  \emph{Errors:} 400 \code{MISSING_FIELDS}, \code{INVALID_FIELD};
  429 \code{BUDGET_EXCEEDED}; 502 \code{MODEL_ERROR}. \\
\addlinespace
\apiroute{GET}{/api/llm/models}
admin & The models the gateway offers, as the assistant service reports them: \code{providers}
  and \code{models}, each with \code{id}, \code{name}, \code{kind} (\code{chat} or
  \code{embedding}), \code{contextLength}, \code{maxOutputTokens} and \code{pricing}. \\
\end{xltabular}

\subsection{MCP servers}\label{sec:api:mcp}

The registry of external tool servers whose tools the chat assistant may call, for
administrators (\file{assistant/api/mcp_views.py}; \cref{ch:chat}). A server's header values
are credentials: they are never returned, only their names (\code{headerNames}).

\begin{xltabular}{\textwidth}{@{}L{19mm}Y@{}}
\caption{Routes: MCP servers}\label{tab:api:mcp}\\
\toprule \hd{Who} & \hd{What the route does} \\ \midrule \endfirsthead
\toprule \hd{Who} & \hd{What the route does} \\ \midrule \endhead
\bottomrule \endlastfoot
\apiroute{GET POST}{/api/mcp/servers}
admin & \code{POST}: \code{name} (2 to 31 characters: lower-case letters, digits and the
  underscore, a letter first), \code{url}, and optionally \code{transport} (\code{http} or
  \code{sse}), \code{headers} (at most 10), \code{allowedTools} (empty means every tool of the
  server), \code{description}, \code{enabled}. The URL must be https on the standard port to a
  public DNS name --- no address, no container or local name, no credentials in it. A new
  server starts disabled. \emph{Errors:} 400 \code{MISSING_FIELDS}, \code{INVALID_FIELD}; 409
  \code{DUPLICATE_NAME}. \\
\addlinespace
\apiroute{GET PATCH DELETE}{/api/mcp/servers/{id}}
admin & \code{PATCH}: any subset of the same fields. Headers are merged: a value of
  \code{null} removes that header and the others are kept. A server whose stored URL breaks
  the rule cannot be enabled. \emph{Errors:} 400 \code{INVALID_FIELD}; 409
  \code{DUPLICATE_NAME}. \\
\addlinespace
\apiroute{POST}{/api/mcp/servers/{id}/test}
admin & The assistant service connects to the server and lists its tools. The answer is always
  200 with the server and \code{ok}: on success \code{tools} holds the listing, on failure
  \code{lastError} says why. Audited either way. \emph{Errors:} 400 \code{INVALID_FIELD} (the
  stored URL breaks the rule). \\
\end{xltabular}

\subsection{Jobs}\label{sec:api:jobs}

The monitor of the worker's queue (\file{jobs/api/views.py}; \cref{ch:arch}). A user sees the
jobs they caused, an administrator all.

\begin{xltabular}{\textwidth}{@{}L{19mm}Y@{}}
\caption{Routes: jobs}\label{tab:api:jobs}\\
\toprule \hd{Who} & \hd{What the route does} \\ \midrule \endfirsthead
\toprule \hd{Who} & \hd{What the route does} \\ \midrule \endhead
\bottomrule \endlastfoot
\apiroute{GET}{/api/jobs}
owner & \code{jobs}, newest first; filters \code{status}, \code{kind}, \code{subject};
  \code{limit}. A job about an act carries the act's name as \code{subjectTitle}.
  \emph{Errors:} 400 \code{INVALID_QUERY}. \\
\addlinespace
\apiroute{GET}{/api/jobs/counts}
owner & How many jobs the caller can see in each status: \code{queued}, \code{running},
  \code{done}, \code{failed}, \code{cancelled}. \\
\addlinespace
\apiroute{GET}{/api/jobs/{id}}
owner & One job with its \code{steps}: for each, its name, status, duration, error, the prompt
  version it ran on, and what its model calls used (\code{usage}). \\
\addlinespace
\apiroute{POST}{/api/jobs/{id}/cancel}
owner & A queued job is cancelled at once; a running one gets \code{cancelRequested} and stops
  at its next check. Audited. \emph{Errors:} 409 \code{NOT_CANCELLABLE} (already finished). \\
\end{xltabular}

\subsection{Usage}\label{sec:api:usage}

What the model calls cost, per day and per person (\file{assistant/api/usage_views.py};
\cref{ch:models}). Everyone reads their own usage. An administrator widens the view to every
account with \code{scope=all} or looks at one account with \code{userId}; a user who asks for
either gets 403 \code{FORBIDDEN}.

\begin{xltabular}{\textwidth}{@{}L{19mm}Y@{}}
\caption{Routes: usage}\label{tab:api:usage}\\
\toprule \hd{Who} & \hd{What the route does} \\ \midrule \endfirsthead
\toprule \hd{Who} & \hd{What the route does} \\ \midrule \endhead
\bottomrule \endlastfoot
\apiroute{GET}{/api/usage}
own; admin: wider & The last \code{days} days (1 to 366, 30 when absent) in UTC, a row for
  every day: \code{days}, \code{byKind}, \code{byModel}, \code{totals}, \code{from},
  \code{to}; \code{byUser} in an administrator's wide view; and \code{budget} --- today's
  \code{limit}, \code{used} and \code{remaining} of the account looked at (\code{remaining} is
  \code{null} when unlimited). \emph{Errors:} 400 \code{INVALID_QUERY}. \\
\addlinespace
\apiroute{GET}{/api/usage/calls}
own; admin: wider & The log of model calls, newest first: \code{calls} and \code{total};
  filter \code{kind} (\code{chat}, \code{job_step}, \code{embed}, \code{prompt_test}). Each
  call carries its model, tokens (input, read from the cache, written to it, output,
  reasoning), \code{costUsd}, duration and outcome. \emph{Errors:} 400 \code{INVALID_QUERY}. \\
\addlinespace
\apiroute{GET}{/api/usage/calls/{id}}
own; admin: all & One call with its \code{toolCalls} and, for a chat turn, the
  \code{transcript}: what the user asked and what the assistant answered. \\
\end{xltabular}

\subsection{Audit}\label{sec:api:audit}

The trail of who did what (\file{audit/api/views.py}; \cref{ch:security}).

\begin{xltabular}{\textwidth}{@{}L{19mm}Y@{}}
\caption{Routes: audit}\label{tab:api:audit}\\
\toprule \hd{Who} & \hd{What the route does} \\ \midrule \endfirsthead
\toprule \hd{Who} & \hd{What the route does} \\ \midrule \endhead
\bottomrule \endlastfoot
\apiroute{GET}{/api/audit}
admin & \code{entries}, newest first, each with \code{id}, \code{userId}, \code{userName},
  \code{action}, \code{targetType}, \code{targetId}, \code{details}, \code{ip},
  \code{createdAt}; and \code{next}. Filters: \code{action}, a prefix (\code{users.} for every
  change of an account), and \code{user}, one account's id. \code{next} holds \code{before} and
  \code{beforeId} for the following page, to be passed back as they are, and is \code{null}
  after the last page. \emph{Errors:} 400 \code{INVALID_QUERY}. \\
\end{xltabular}


% ============================================================
\section{The assistant service's public routes}\label{sec:api:assistant}
% ============================================================

The assistant service answers three routes to the outside, all under
\code{/api/assistant/chat} (in the assistant service, \file{routes/chat.js}; mounted in
\file{app.js}). They take the same token as the API: the service resolves it by calling
\route{GET}{/api/auth/me} and remembers the answer for 15 seconds. What happens inside a turn
--- the prompt, the tools, the model --- is the subject of \cref{ch:chat}.

\begin{xltabular}{\textwidth}{@{}L{19mm}Y@{}}
\caption{Routes: the assistant service, public}\label{tab:api:assistant}\\
\toprule \hd{Who} & \hd{What the route does} \\ \midrule \endfirsthead
\toprule \hd{Who} & \hd{What the route does} \\ \midrule \endhead
\bottomrule \endlastfoot
\apiroute{POST}{/api/assistant/chat}
signed in, own thread & One chat turn. The body carries one new message of the user; the
  conversation before it is read from the API. The answer is a stream
  (\code{text/event-stream}) in the format of the chat runtime's UI messages. \\
\addlinespace
\apiroute{POST}{/api/assistant/chat/{threadId}/stop}
signed in, own thread & Stops the caller's live turn in that thread. Answers \code{stopped},
  \code{true} or \code{false}; the partial reply is stored with the status \code{stopped}. A
  thread of another user reads as nothing to stop. \\
\addlinespace
\apiroute{GET}{/api/assistant/chat/{threadId}/status}
signed in, own thread & Whether the thread has a live turn: \code{running} and
  \code{startedAt}. The console polls it after a page is reloaded in the middle of an answer.
  Never cached. \\
\end{xltabular}

\paragraph{The body of a turn.} \code{readTurnBody} accepts this shape and rebuilds it field
by field, so that nothing unexpected reaches the API or the model:

\begin{codeblock}
{"id": "<thread id>",
 "trigger": "submit-message" | "regenerate-message",
 "message": {"id": "<message id>", "role": "user",
             "parts": [{"type": "text", "text": "..."},
                       {"type": "data-document", "data": {"id": "<document id>", "title": "...",
                        "kind": "...", "charCount": 1234, "extractionStatus": "done"}}],
             "metadata": {...}},
 "parentId": "<message id>" | null,
 "locale": "lt" | "en"}
\end{codeblock}

The thread id must be a uuid; \code{trigger} is \code{submit-message} when absent; the
message's optional \code{metadata} is kept only while it is an object of at most 2\,000
characters as JSON. A message has 1 to 30 parts (\code{MAX_PARTS}); the text parts together hold at most
\code{MAX_MESSAGE_CHARS}~=~60\,000 characters. A \code{data-document} part is a marker for a
document uploaded beforehand through \route{POST}{/api/documents}; a file itself is never part
of a message. The message id is chosen by the client, from letters, digits and
\code{_ . : -}, at most 128 characters.

\paragraph{Before the stream.} The body parser comes first: 400 \code{INVALID_JSON} for a
body that is not JSON, 413 \code{PAYLOAD_TOO_LARGE} beyond 512~kB. Then a turn passes its
gates in a fixed order, and a failure at any of them is an ordinary JSON error in the
assistant service's envelope: the token (401 \code{AUTH_REQUIRED} without one,
\code{SESSION_INVALID} for a dead one); a configured model (503 \code{NOT_CONFIGURED}); the
rate limit (429 \code{RATE_LIMITED}); the body (400 \code{INVALID_BODY}, \code{INVALID_THREAD},
\code{INVALID_TRIGGER}, \code{INVALID_MESSAGE}, \code{INVALID_PARENT}, \code{INVALID_PART},
\code{INPUT_TOO_LARGE}, \code{EMPTY_MESSAGE}); the daily budget (429 \code{BUDGET_EXCEEDED};
a budget check that fails lets the turn through); an active version of \prompt{chat.assistant} (503
\code{PROMPT_NOT_CONFIGURED}, or \code{PROMPT_UNAVAILABLE} when it cannot be loaded); and the
thread, fetched from the API as the user (404 \code{THREAD_NOT_FOUND}). When the API cannot be
reached the answer is 502 \code{DJANGO_UNREACHABLE}; any other refusal of the API passes
through with the API's own status and code. A service that is shutting down answers 503
\code{SHUTTING_DOWN}.

\paragraph{Inside the stream.} Once the stream has begun the HTTP status is 200 whatever
happens, and a failure arrives as an error chunk of the stream carrying a code:
\code{MODEL_TIMEOUT} (a turn runs at most \code{CHAT_TIMEOUT_MS}~=~900 seconds),
\code{AI_BUDGET_EXHAUSTED} (the gateway's own budget is spent), \code{MODEL_UNAVAILABLE} (a
passing failure: trying again later can help), \code{MODEL_REFUSED},
\code{MODEL_CONTEXT_EXCEEDED}, \code{MODEL_CONTENT_FILTERED}, \code{MODEL_ACCESS_DENIED}, or
\code{STREAM_FAILED} for anything else. The turn does not depend on the connection: if the
browser leaves, the answer is still completed and stored, and the reply's metadata records how
the turn ended when it did not end well (\code{stopped}, \code{interrupted}, \code{timeout},
\code{budget}, \code{error}, or the kind of closing line the service had to add:
\code{step-limit}, \code{length}, \code{content-filter}, \code{empty}).


% ============================================================
\section{The internal doors}\label{sec:api:internal}
% ============================================================

The API and the assistant service need each other: the worker and the API ask the assistant
service for every model call, and the assistant service asks the API for the chat's prompt and
budget and hands it the replies. These calls go through two \term{internal doors}, the paths
under \code{/internal/} of each service (\cref{fig:api:doors}). Neither door can be reached from
outside, for two independent reasons:

\begin{itemize}
  \item \emph{No path leads there.} The ingress forwards \code{/api/assistant/*} to the
    assistant service and \code{/api/*} to the API, and nothing else to either. A request for
    \code{/internal/...} from outside falls into the third family and is answered by the web
    server of the console. The two services publish no port of their own.
  \item \emph{The shared secret.} Every internal route demands the header
    \code{X-Internal-Secret}, which must equal the value both services receive as
    \code{ASSISTANT_INTERNAL_SECRET} when the installation starts. The comparison takes
    constant time, and a service whose secret is empty refuses every internal request: the
    door fails closed. The refusal is 403 \code{INTERNAL_ONLY} (\code{require_internal} in
    \file{assistant/api/internal_views.py}; in the assistant service, \code{requireInternal}
    in \file{middleware/internal.js}, which runs before the body is even parsed).
\end{itemize}

\begin{figure}[H]
\centering
\begin{tikzpicture}[node distance=8mm and 50mm]
  \node[svc, minimum width=34mm, minimum height=22mm] (django) {the API\\and the worker\\(Django)};
  \node[svc, minimum width=34mm, minimum height=22mm, right=of django] (assistant) {assistant\\service};
  \node[ext, right=14mm of assistant] (gw) {gateway};
  \draw[call] ([yshift=7mm]django.east) -- ([yshift=7mm]assistant.west)
    node[lbl, midway, above=1pt] {\code{/internal/llm/*}, \code{/internal/mcp/*}};
  \draw[call] (assistant.west) -- (django.east)
    node[lbl, midway] {\code{/internal/assistant/*}};
  \draw[flow] ([yshift=-7mm]assistant.west) -- ([yshift=-7mm]django.east)
    node[lbl, midway, below=1pt] {\code{/api/*} with the user's token};
  \draw[call] (assistant) -- (gw);
  \begin{scope}[on background layer]
    \node[lane, fit=(django)(assistant), inner ysep=9mm] (lane) {};
  \end{scope}
  \node[lanelabel, anchor=south west] at (lane.north west) {isolated network};
  \node[note, anchor=north west] at ([yshift=-1mm]lane.south west) {between the two services --- dashed: the header \code{X-Internal-Secret} is required; solid: the public API, under the user's own rights};
\end{tikzpicture}
\caption{The internal doors. The API and its worker call the assistant service for every
model call; the assistant service calls back for what only the API knows. Whatever concerns a
user's own data the assistant service fetches through the public routes with that user's
token, so the tenancy rule is applied in one place.}\label{fig:api:doors}
\end{figure}

\subsection{The door of the API}\label{sec:api:door-django}

Six routes under \code{/internal/assistant/} serve the assistant service
(\file{assistant/api/internal_views.py}). Their field names are snake\_case.

\begin{xltabular}{\textwidth}{@{}L{19mm}Y@{}}
\caption{Routes: the internal door of the API}\label{tab:api:door-django}\\
\toprule \hd{Who} & \hd{What the route does} \\ \midrule \endfirsthead
\toprule \hd{Who} & \hd{What the route does} \\ \midrule \endhead
\bottomrule \endlastfoot
\apiroute{GET}{/internal/assistant/prompt}
the secret & The active version of the prompt named by \code{slug}, with its text, model and
  settings. 404 \code{PROMPT_NOT_CONFIGURED} when the prompt has no active version --- the chat
  then answers 503 rather than run on instructions nobody can see. \\
\addlinespace
\apiroute{POST}{/internal/assistant/turn-log}
the secret & One row of the usage log for a chat turn: \code{user_id}, \code{thread_id},
  \code{model}, \code{prompt_version_id}, the token counts (\code{input_tokens},
  \code{cached_input_tokens}, \code{cache_write_tokens}, \code{output_tokens},
  \code{reasoning_tokens}), \code{duration_ms}, \code{tool_calls}, \code{outcome} and the
  \code{transcript} (the user's text cut to 20\,000 characters, the answer to 100\,000).
  The API computes the cost. Answers 201 with \code{id} and the user's \code{budget}. \\
\addlinespace
\apiroute{GET}{/internal/assistant/budget}
the secret & Today's budget of the user \code{user_id}: \code{limit}, \code{used},
  \code{remaining} (\code{null} when unlimited), \code{unit}. Asked before every turn. \\
\addlinespace
\apiroute{POST}{/internal/assistant/threads/{id}/messages}
the secret & Stores one reply of the assistant: \code{user_id}, \code{id}, \code{parent_id},
  \code{format}, \code{content}. This is the only writer of assistant messages, so a browser
  can never forge the output of a tool. The thread must belong to \code{user_id} (404
  otherwise). 201 when the message is new, 200 when it replaces one. \\
\addlinespace
\apiroute{POST}{/internal/assistant/threads/{id}/title}
the secret & The title a model wrote for a thread: \code{user_id}, \code{title}. It is applied
  only while the thread still carries its automatic title, so a title the user set stays.
  Answers \code{ok} and \code{applied}. \\
\addlinespace
\apiroute{GET}{/internal/assistant/mcp-servers}
the secret & The enabled MCP servers \emph{with} their header values --- the only place those
  credentials leave the API. A server whose stored URL breaks the address rule is left out. \\
\end{xltabular}

\subsection{The door of the assistant service}\label{sec:api:door-assistant}

Four routes under \code{/internal/} serve the API and its worker (in the assistant service,
\file{routes/internal.js} and \file{routes/internalMcp.js}); the client on the API's side is
\file{assistant/gateway.py}. They are the system's only way to a model, and \cref{ch:models}
describes them in full. A fifth route, \code{/health}, needs no secret and is not forwarded by
the ingress either.

\begin{xltabular}{\textwidth}{@{}L{19mm}Y@{}}
\caption{Routes: the internal door of the assistant service}\label{tab:api:door-assistant}\\
\toprule \hd{Who} & \hd{What the route does} \\ \midrule \endfirsthead
\toprule \hd{Who} & \hd{What the route does} \\ \midrule \endhead
\bottomrule \endlastfoot
\apiroute{POST}{/internal/llm/generate}
the secret & One model call. Body \code{prompt}, and optionally \code{system}, \code{shared},
  \code{model}, \code{schema}, \code{schemaName}, \code{temperature}, \code{maxOutputTokens}.
  Answers \code{model}, \code{text} or \code{object} (with a schema, the answer checked against
  it), \code{finishReason}, \code{usage}, \code{durationMs}. A call is given up after
  \code{GENERATE_TIMEOUT_MS} (900 seconds by default). \emph{Errors:} 400 for a malformed
  request; 402 \code{AI_BUDGET_EXHAUSTED} when the gateway's own budget is spent; otherwise
  \code{GENERATION_FAILED} with the gateway's status, or 502 --- and 429 with
  \code{retryAfterS} for a rate limit. \\
\addlinespace
\apiroute{POST}{/internal/llm/embed}
the secret & Embeddings for \code{texts}, at most 256 in a call; optionally \code{model}.
  Answers \code{model}, \code{embeddings}, \code{dimensions}, \code{usage},
  \code{durationMs}. \emph{Errors:} 400 \code{INVALID_TEXTS}, \code{TOO_MANY_TEXTS},
  \code{UNKNOWN_MODEL}; \code{EMBEDDING_FAILED}; 402 \code{AI_BUDGET_EXHAUSTED}. \\
\addlinespace
\apiroute{GET}{/internal/llm/models}
the secret & \code{providers} and \code{models}: the gateway's list of models with their
  prices, newest first, kept for ten minutes. \\
\addlinespace
\apiroute{POST}{/internal/mcp/list-tools}
the secret & Connects to the MCP server described by \code{url}, \code{transport} and
  \code{headers} and lists its tools --- the connection test of
  \route{POST}{/api/mcp/servers/{id}/test}. \emph{Errors:} 400 \code{INVALID_URL},
  \code{INVALID_TRANSPORT}, \code{INVALID_HEADERS}; 502 \code{MCP_UNREACHABLE}, with the
  credentials cut out of the message. \\
\addlinespace
\apiroute{GET}{/health}
inside the network & A probe of liveness and configuration (the compose file sets no health check):
  \code{status}, \code{model} (\code{configured} or \code{missing-key}), \code{providers},
  \code{activeTurns}. \\
\end{xltabular}


% ============================================================
\section{What a client has to know}\label{sec:api:client}
% ============================================================

Two behaviours of the API are not visible in a single route: work that outlasts the request
that started it, and answers that are files.

\subsection{Long work: start, poll, decide, cancel}\label{sec:api:polling}

No route keeps a client waiting for an ingestion, an extraction, a report or a work product.
The route records what is wanted, queues a job for the worker and answers at once; the client
then asks again until the work has ended. There is no push channel: the chat's stream is the
only stream in the system. \Cref{fig:api:polling} shows the sequence for a report.

\begin{figure}[H]
\centering
\begin{tikzpicture}[node distance=5mm and 6mm]
  \node[svc] (s1) {\code{queued}};
  \node[svc, right=of s1] (s2) {\code{running}};
  \node[man, right=of s2] (s3) {\code{scope_review}};
  \node[svc, right=of s3] (s4) {\code{queued}};
  \node[svc, right=of s4] (s5) {\code{running}};
  \node[svc, right=of s5] (s6) {\code{done}};
  \draw[flow] (s1) -- (s2); \draw[flow] (s2) -- (s3); \draw[flow] (s3) -- (s4);
  \draw[flow] (s4) -- (s5); \draw[flow] (s5) -- (s6);
  \draw[opt] (s2.south) -- ++(0,-5mm) -| (s5.south) node[lbl, pos=.25] {without \code{review}};
  \node[data, above=11mm of s1] (c1) {\code{POST /api/reports}};
  \node[data, above=11mm of s3] (c3) {\code{POST .../scope}};
  \node[data, above=11mm of s6] (c6) {\code{GET .../export}};
  \node[data] (c2) at ($(c1)!.5!(c3)$) {\code{GET} \ldots};
  \node[data] (c5) at ($(c3)!.5!(c6)$) {\code{GET} \ldots};
  \draw[call] (c1) -- (s1) node[lbl, midway] {201};
  \draw[call] (c3) -- (s3) node[lbl, midway] {approve};
  \draw[call] (c6) -- (s6) node[lbl, midway] {file};
  \draw[call] (c2) -- (s2.north);
  \draw[call] (c5) -- (s5.north);
  \node[lanelabel, anchor=south west] at ([yshift=1mm]c1.north west) {the client's calls};
  \node[lanelabel, anchor=north west] at ([yshift=-7mm]s1.south west) {the report's \code{status}; a report can also end \code{failed} or \code{cancelled}};
\end{tikzpicture}
\actorlegend
\caption{A report as a client sees it. The start answers at once; the client polls
\code{GET /api/reports/{id}} for \code{status} and \code{progress}; with \code{review} the
report stops and waits for a person at the scope route; a finished report can be
exported.}\label{fig:api:polling}
\end{figure}

\paragraph{What the start answers.} Starting a report or a work product answers 201 with the
new row, already carrying \code{jobId}, \code{jobStatus} and \code{progress}. The routes that
change the corpus --- refresh, upload, obligation extraction, the judge, re-embedding,
refresh-all --- answer 202 with the job itself. A document upload is the exception in the other
direction: the text is extracted before the answer, and only the indexing that follows is a
job, which the answer does not name; it is found with \route{GET}{/api/jobs} and
\code{subject=document:<id>}.

\paragraph{Twice the same request.} The API refuses or folds a second start rather than run
the same work twice. A second report of a company, or a second product of a kind, answers 409
\code{ALREADY_RUNNING}. A second refresh, extraction or judge run on the same act answers 202
with the job that is already queued or running, so a client can treat the two answers alike.

\paragraph{What to poll.} Three resources show the state of long work
(\cref{tab:api:polling}). \code{progress} is the job's own report of where it is --- the step,
a percentage and a message. A job additionally lists its steps, each with what its model calls
used, once the step has ended.

\begin{table}[H]
\caption{The resources a client polls, and the values of their status}\label{tab:api:polling}
\begin{tabularx}{\textwidth}{@{}L{38mm}Y@{}}
\toprule
\hd{Resource} & \hd{\code{status} is one of} \\
\midrule
\code{GET /api/jobs/{id}} & \code{queued}, \code{running}, \code{done}, \code{failed}, \code{cancelled}; with \code{progress}, \code{result}, \code{error}, \code{cancelRequested}, \code{attempts} \\
\code{GET /api/reports/{id}} & \code{queued}, \code{running}, \code{scope_review}, \code{done}, \code{failed}, \code{cancelled}; with \code{jobStatus}, \code{progress}, \code{error} \\
\code{GET /api/programme/{id}} & \code{queued}, \code{running}, \code{draft}, \code{approved}, \code{failed}, \code{cancelled}; with \code{jobStatus}, \code{progress}, \code{error} \\
\bottomrule
\end{tabularx}
\end{table}

The list routes carry the same fields, so a client that shows many rows polls
\route{GET}{/api/reports}, \route{GET}{/api/companies/{id}/programme} or \route{GET}{/api/jobs}
instead of each row. An act's detail carries \code{latestJob} --- the ingestion or extraction
at work on it now or, when none is, the most recent one. For the chat, the status route of the
assistant service plays the same part after a reload (\cref{sec:api:assistant}).

\paragraph{Where a person is asked.} A report started with \code{review: true} does not run
through: after deciding which acts apply it stops in \code{scope_review} and no job is at work
on it. It continues only when \route{POST}{/api/reports/{id}/scope} is called with
\code{approved: true}, which queues the job again. The same route without \code{approved} can
be called any number of times to price a selection; it stores nothing. A work product ends its
generation as a \code{draft}, and becomes \code{approved} only by a person's
\route{PATCH}{/api/programme/{id}}.

\paragraph{Cancelling.} Each of the three resources has a cancel route. A job, report or
product that is still queued is cancelled at once; one that is running gets a flag that the
worker honours at its next check, so the status changes a little later; a report waiting for
its scope is cancelled at once. Anything already ended answers 409 \code{NOT_CANCELLABLE}.
Deleting a report or a product in work is refused with 409 \code{STILL_RUNNING}: it is
cancelled first.

\paragraph{Money before work.} A client can ask before it commits the user's budget:
\route{GET}{/api/reports/estimate} gives the upper bound of a report's cost together with the
budget left today, the scope route prices the second part for a selection of acts, and
\route{GET}{/api/usage} carries today's \code{budget}. A started run is never stopped halfway
for lack of budget; the budget only decides whether it may start.

\subsection{Files}\label{sec:api:files}

Three routes answer a file instead of JSON, built at the moment of the request
(\cref{tab:api:exports}). Each sets \code{Content-Disposition: attachment} with a file name and
\code{Cache-Control: no-store}, and each is written to the audit trail. In the file name the
Lithuanian letters of the company's name are folded to plain Latin ones, every other
run of characters outside letters, digits, dot, underscore and hyphen becomes one hyphen, and
the name before its extension is cut to 80 characters (for a report, before the prefix
\code{legalops-} is added). An error on these routes is still the JSON envelope, so a client
checks the status before it treats the body as a file.

\begin{table}[H]
\caption{The routes that answer a file}\label{tab:api:exports}
\begin{tabularx}{\textwidth}{@{}L{47mm}L{24mm}Y@{}}
\toprule
\hd{Route} & \hd{Formats} & \hd{File name; condition} \\
\midrule
\code{GET /api/reports/{id}/export} & \code{pdf}, \code{xlsx} & \code{legalops-<company>-<date>}; only a report that is \code{done} \\
\code{GET /api/companies/{id}/actions/export} & \code{xlsx} & \code{actions-<company>-<date>}; always \\
\code{GET /api/programme/{id}/export} & \code{pdf}; \code{docx} (policy); \code{xlsx} (risks) & \code{<kind>-<company>-v<version>}; a \code{draft} or an \code{approved} product \\
\bottomrule
\end{tabularx}
\end{table}

The organisation chart of a governance product (\route{GET}{/api/programme/{id}/chart}) is the
one drawing the API serves: an SVG image sent inline, not as an attachment. Files travel into
the system through two routes only, both \code{multipart/form-data} with the part \code{file}:
\route{POST}{/api/documents} for a company's document and
\route{POST}{/api/corpus/documents/{id}/upload} for a new version of an act. Every other
request body is JSON.
